% ================================================================
\section{Glossary}
\label{app:glossary}
% ================================================================

Terms are grouped by subject and listed alphabetically within each group. Each entry gives a
plain-language definition and, where useful, a reference for further reading. Introductory
textbooks for each area are: market microstructure~\citep{harris2003,ohara1995,hasbrouck2007};
futures~\citep{hull}; statistical learning~\citep{hastie2009}; deep learning~\citep{goodfellow2016};
point processes~\citep{daley2003}.

\newcommand{\gl}[1]{\item[#1]}

\subsection{Markets and trading}
\begin{description}[leftmargin=1.2em,style=sameline,itemsep=1pt,font=\normalfont\bfseries]
\gl{Adverse selection} The tendency of a passive order to be filled exactly when the price is
  about to move against it, because the counterparty often knows more~\citep{glosten1985}.
  Section~\ref{sec:adverse}.
\gl{Aggressive / passive} An aggressive order trades immediately against waiting orders; a passive
  order waits in the book for someone else to trade with it.
\gl{Alpha} The return of a strategy or signal beyond an appropriate benchmark, after the relevant risk
  and cost adjustments; in this survey, usually the part of a short-horizon forecast that is profitable
  after costs.
\gl{Arrival price} The mid-price at the moment an order is decided, used as the reference price for
  its execution cost.
\gl{Backtest} Running a trading strategy on historical data to estimate how it would have done.
\gl{Basis point (bp)} One hundredth of a percent, 0.01\%. A parameter of 300\,bp is 3\%.
\gl{Bid and ask} A waiting buy order is a bid and a waiting sell order an ask (or offer); the best
  bid is the highest bid price and the best ask the lowest ask price.
\gl{Capacity} How much a strategy can trade before its own price impact erases its profit.
\gl{Clearing fees} Fees paid to the clearing house, the institution that settles trades and
  guarantees both sides.
\gl{Counterparty} The trader on the other side of a trade.
\gl{Cross-sectional} Across many instruments at the same moment, rather than over time for one.
\gl{Depth} The total size waiting in the book near the best prices; a \emph{deep} book has a lot.
\gl{Drift} The average direction in which the price moves over a period, independent of one's
  own trades.
\gl{Edge} Expected profit per trade, before or after costs as stated.
\gl{Equity / fixed income} Stocks / bonds and other interest-rate products.
\gl{ETF} Exchange-traded fund: a share, traded like a stock, that holds a basket of assets.
\gl{Fill} The execution of an order, in full or in part.
\gl{Futures contract} A standardised exchange-traded agreement to buy or sell an asset at a set
  future date and price~\citep{hull}. ES, NQ and ZN are futures on the S\&P~500 index, the
  Nasdaq-100 index and the 10-year US Treasury note.
\gl{Iceberg / hidden order} An order that shows only part (iceberg) or none (hidden) of its size.
\gl{Implementation cost / slippage} The difference between the price obtained and a reference
  price such as the mid at the moment of decision~\citep{harris2003}.
\gl{Index constituents} The individual stocks that make up an index such as the S\&P~500.
\gl{Informed trader} A trader who knows something about the next price move that others do not.
\gl{Instrument} One specific tradable product, for example the March ES contract.
\gl{Inventory / position} The net quantity held after one's fills, long (positive) or short
  (negative); it carries risk until offset.
\gl{Large-tick asset} An instrument whose tick is large relative to its typical price moves, so the
  spread is almost always one tick and queues are long.
\gl{Lifting the offer / hitting the bid} Buying aggressively at the ask / selling aggressively at
  the bid.
\gl{Limit order} An order to buy or sell up to a given quantity at a given price or better; what
  cannot trade at once waits in the book.
\gl{Liquidity} How easily one can trade without moving the price. Resting limit orders supply it;
  aggressive orders consume it.
\gl{Market impact} How much one's own trading moves the price against oneself~\citep{bouchaud2018}.
\gl{Market making} Continuously posting both a bid and an ask in order to earn the spread, while
  managing inventory~\citep{avellaneda2008,cartea2015}. A posted bid or ask is a \emph{quote}.
\gl{Market order} An order to buy or sell a quantity now at whatever price is available; it trades
  at once against waiting orders and never waits itself.
\gl{Mark-out} The gain or loss of a fill measured against the mid a fixed time later.
  Equation~\eqref{eq:markout}.
\gl{Mass cancel} A single command cancelling all of a trader's open orders.
\gl{Metaorder} A large order executed as a sequence of smaller orders over time.
\gl{Micro-price} An estimate of the ``fair'' price that adjusts the mid toward the side with the
  smaller queue~\citep{stoikov2018}.
\gl{Mid-price} The average of the best bid and the best ask.
\gl{Optimal execution} Splitting a large order over time to minimise its total cost and
  risk~\citep{almgren2001}.
\gl{Paper trading} Running a strategy on live data with simulated, not real, orders.
\gl{P\&L} Profit and loss.
\gl{Pre-trade checks} Automatic checks that block an order breaking size, price or position limits
  before it is sent.
\gl{Price point} The contract's own price unit. One ES point is four ticks, worth \$50 per
  contract.
\gl{Queue position} The volume waiting ahead of an order in the queue at its price.
\gl{Quote skew} Shifting both a market maker's bid and ask in the direction of expected price movement or to reduce inventory.
\gl{Rebate} A payment some venues make to passive orders for supplying liquidity, funded by fees on
  aggressive orders.
\gl{Regime} A stretch of time with stable market behaviour, such as calm or volatile.
\gl{Return} The relative change in value over a period.
\gl{Session} One period during which the exchange is open, usually one trading day.
\gl{Sharpe ratio} Average return divided by the standard deviation of return: profit per unit of
  risk~\citep{sharpe1966,sharpe1994}.
\gl{Spoofing / layering} Placing orders one intends to cancel before they trade---one large order
  (spoofing) or several stacked across prices (layering)---to create a false impression of supply
  or demand. Illegal in US futures markets~\citep{cftc2013spoofing}.
\gl{Spread} The difference between the best ask and the best bid.
\gl{Square-root impact law} The empirical regularity that the price impact of a metaorder grows roughly with the square root of its size~\citep{bouchaud2018}.
\gl{Stress testing} Checking a system or strategy under extreme or rare scenarios.
\gl{The touch} The current best bid and best ask.
\gl{Tick} The smallest allowed price step of an instrument.
\gl{Venue} A trading platform or exchange.
\gl{Volatility} How much the price fluctuates, usually the standard deviation of returns.
\gl{VWAP} Volume-weighted average price: the average trade price over an interval, each trade
  weighted by its size~\citep{harris2003}.
\gl{Zero-intelligence model} A model in which orders arrive at random with no strategy; the matching rules alone generate the book's behaviour~\citep{smith2003}.
\end{description}

\subsection{Exchanges and market data}
\begin{description}[leftmargin=1.2em,style=sameline,itemsep=1pt,font=\normalfont\bfseries]
\gl{CME Group} The Chicago Mercantile Exchange group, the main US futures exchange operator; its
  electronic platform is CME Globex.
\gl{Event clock / calendar clock} Sampling the market after a fixed number of events / a fixed number of seconds.
\gl{Feed / feed handler} An exchange's real-time stream of market-data messages / the software
  that decodes the feed and maintains the book.
\gl{iLink 3} CME's order-entry protocol, used to send, modify and cancel orders~\citep{cme_ilink3}.
\gl{ITCH} Nasdaq's market-by-order data protocol~\citep{itch50}; Section~\ref{sec:intro}.
\gl{L1, L2, L3} Feed detail levels: best bid and ask only; total size at each price level (market
  by price, MBP); every individual order (market by order, MBO).
\gl{LOBSTER} A service that reconstructs Nasdaq order books from ITCH data for
  research~\citep{lobster}.
\gl{Matching engine} The exchange computer that holds the book and matches buy and sell orders.
\gl{MDP 3.0} CME's Market Data Platform protocol: binary messages over UDP
  multicast~\citep{cme_mdp3}.
\gl{Order acknowledgement} The exchange's confirmation that it has accepted an order.
\gl{Packet capture (PCAP)} A file of network packets recorded exactly as they arrived, with timestamps; replaying
  it reproduces a market-data feed.
\gl{Price--time priority (FIFO)} The matching rule in which better prices trade first and, at the same price, earlier
  orders trade first: first in, first out.
\gl{Snapshot recovery} A periodic full-book message that lets a feed handler resynchronise after
  lost packets.
\end{description}

\subsection{Statistics and evaluation}
\begin{description}[leftmargin=1.2em,style=sameline,itemsep=1pt,font=\normalfont\bfseries]
\gl{Ablation} Removing or adding one component of a model to measure what it contributes.
\gl{AUC} Area under the ROC (receiver operating characteristic) curve: the probability that a randomly chosen positive case receives a
  higher score than a randomly chosen negative one~\citep{fawcett2006}.
\gl{Balanced accuracy} The average of the recall (true-positive rate) of each class.
\gl{Brier score} The mean squared difference between predicted probabilities and 0/1
  outcomes~\citep{brier1950}.
\gl{Calibration} Agreement between predicted probabilities and observed frequencies: events given
  70\% happen about 70\% of the time~\citep{guo2017}.
\gl{Calibration error} A summary of how far predicted probabilities are from observed frequencies~\citep{guo2017}.
\gl{Coefficient of determination ($R^2$)} The fraction of the variance of a target explained by a fitted model.
\gl{Concordance (C-index)} The fraction of pairs whose predicted order of event times matches the
  actual order.
\gl{Confidence interval} A range that contains the true value with a stated probability over
  repeated experiments; it can be computed by resampling~\citep{efron1993}.
\gl{Confounding} Mixing of two effects so that the cause of an observed difference cannot be
  isolated.
\gl{Counterfactual} What would have happened under an action that was not actually taken.
\gl{Cross-validation (random folds)} Splitting data into several parts, training on all but one and testing on the
  remaining one, in rotation~\citep{hastie2009}. Inappropriate for time-ordered market data
  (Section~\ref{sec:deep-real}).
\gl{Decile} One tenth of a sorted data set.
\gl{Epps effect} The fall in measured correlation between two assets when it is computed over very short intervals.
\gl{F1 / macro-F1} F1 is the harmonic mean of precision and recall; macro-F1 averages F1 over the
  classes with equal weight.
\gl{False-discovery rate} The expected fraction of claimed discoveries that are
  false~\citep{benjamini1995}.
\gl{Granger causality} A test of whether past values of one series improve forecasts of another beyond its own past.
\gl{In-sample / out-of-sample (OOS)} Measured on the data used to build a model / on data it has
  never seen.
\gl{Interval coverage} The fraction of outcomes that fall inside predicted intervals, compared
  with the intended (\emph{nominal}) fraction.
\gl{Lag} A past value of a variable at a fixed offset in time.
\gl{Lasso} Regression with a penalty on the sum of absolute coefficients, which sets unhelpful coefficients exactly to zero~\citep{tibshirani1996}.
\gl{Lead--lag} A relation in which one instrument's price moves first and a related one follows after a short delay.
\gl{Log-likelihood} The logarithm of the probability a model assigns to the observed data; higher
  means a better fit.
\gl{Log loss} The average of $-\ln$ of the probability a classifier gave to the true class
  (cross-entropy); lower is better~\citep{hastie2009}.
\gl{Model confidence set} The set of models that cannot be statistically distinguished from the best, at a stated confidence level~\citep{hansen2011mcs}.
\gl{MSE / MAE} Mean squared error / mean absolute error.
\gl{Multiple testing} Trying many hypotheses inflates the chance that some look good by
  luck~\citep{white2000,harvey2016}.
\gl{Non-normal returns} Returns whose distribution has fatter tails or more skew than a Gaussian.
\gl{Non-stationary} Having statistical properties that change over time.
\gl{Null hypothesis / significance} The null is the assumption of no effect. A result is
  \emph{significant at 5\%} if it would occur less than 5\% of the time under the null.
\gl{Principal component} A direction in the data along which variance is largest; the first principal component is the best one-number summary of several correlated variables.
\gl{Probability of backtest overfitting} The estimated probability that the configuration that
  looked best in sample performs below the median out of sample~\citep{bailey2014pbo}.
\gl{Proper scoring rule} A score that a forecaster optimises by reporting true beliefs, such as log loss and the Brier score~\citep{gneiting2007}.
\gl{Purging / embargo} Removing training samples whose information overlaps the test period / leaving
  a gap after each test period, both to prevent leakage of future
  information~\citep{lopezdeprado2018}.
\gl{Selective prediction} Acting only on predictions whose confidence exceeds a
  threshold~\citep{geifman2017}.
\gl{Signature plot} Measured volatility plotted against the sampling interval; higher measured volatility at short intervals reveals bid--ask bounce.
\gl{Standard error / annualised} The standard deviation of an estimate / scaled to a one-year
  period.
\gl{Survival model} A model of the time until an event that correctly handles cases where the event
  has not yet happened (\emph{censoring})~\citep{cox1972}.
\gl{Time-rescaling test} A goodness-of-fit test for point-process models: rescaling time by the
  fitted intensity should turn the gaps between events into independent unit-exponential
  variables~\citep{brown2002}.
\gl{Walk-forward evaluation} Train on a past window, test on the next, roll forward and
  repeat~\citep{lopezdeprado2018}.
\end{description}

\subsection{Machine learning}
\begin{description}[leftmargin=1.2em,style=sameline,itemsep=1pt,font=\normalfont\bfseries]
\gl{Attention} A mechanism that weights each input element by its learned relevance to the current
  one~\citep{vaswani2017}.
\gl{Attentive neural process} A model that predicts a distribution for a new input by attending to
  a \emph{context set} of example input--outcome pairs~\citep{kim2019anp}.
\gl{Autoregressive model} A model that generates a sequence one element at a time, each conditioned
  on those before it.
\gl{Axial attention} Attention applied separately along each axis of a multi-dimensional input, such as time and price level.
\gl{Backbone} The main feature-extracting network beneath a task-specific output layer.
\gl{BERT} A Transformer encoder pretrained by predicting masked-out tokens~\citep{devlin2019}.
\gl{CNN / RNN / LSTM} Convolutional networks apply learned sliding filters~\citep{lecun2015};
  recurrent networks carry a hidden state along a sequence; the LSTM is a recurrent network with
  gates for long-range memory~\citep{hochreiter1997}.
\gl{Context length} How many past elements a sequence model sees at once.
\gl{Cross-entropy loss} The negative log-likelihood of the true class under a classifier's predicted
  probabilities; the standard training loss for classification.
\gl{Data augmentation} Enlarging a training set with modified or synthetic examples.
\gl{Diffusion model} A generative model that learns to reverse a gradual process of adding noise to data.
\gl{Embedding} A learned vector representation of an input.
\gl{Encoder} A network that maps an input to a learned vector representation.
\gl{Fine-tuning} Further training a pretrained model on a specific task.
\gl{Foundation model} A large model pretrained on broad data and then adapted to many tasks.
\gl{Generative adversarial network (GAN)} A generator trained against a discriminator that tries to tell generated data from real data.
\gl{Generative model} A model that produces new samples resembling its training data.
\gl{Inception module} A network block that applies several convolutions of different widths in parallel and concatenates the results.
\gl{Inference} Running a trained model to obtain a prediction; \emph{inference latency} is how long
  that takes.
\gl{Logit / logistic regression} The log-odds $\ln\big(\P(A)/(1-\P(A))\big)$ of an event $A$; logistic regression models it as a
  linear function of the inputs~\citep{hastie2009}.
\gl{Meta-model} A model whose inputs are the outputs of other models; here, a learner that combines the outputs of classical order-book models.
\gl{Model capacity} How complex a function a model can represent; roughly, its number of
  parameters.
\gl{Out-of-sample predictions (for stacking)} Predictions of a first-stage model on data it was not fitted on; a second-stage model must be trained on these, not on in-sample predictions, or it learns from the first stage's overfitting~\citep{breiman1996stack}.
\gl{Pretraining} Training a model on a generic task before adapting it to a specific one.
\gl{Quantisation} Using low-precision integers instead of floating point for a network's weights
  and arithmetic, for speed or to fit hardware~\citep{jacob2018}.
\gl{Query, key, value} The three vectors an attention layer computes for each element: the query is
  matched against every key to set the weights, and the output is the weighted average of the values.
\gl{Reinforcement learning} Learning a decision policy by trial and error from rewards.
\gl{Selective state-space model} A state-space model, such as Mamba, whose recurrence parameters depend on the current input~\citep{gu2024mamba}.
\gl{Softmax} A function that turns a vector of scores into probabilities that are positive and sum to one.
\gl{Softplus} The function $\ln(1+\mathrm{e}^{y})$: smooth, always positive, and close to $y$ when $y$ is large; used to keep predicted rates positive.
\gl{Stacking} Training a second-stage model on the predictions of several first-stage models~\citep{wolpert1992,breiman1996stack}.
\gl{State-space model} A model with a hidden state updated by a (learned) linear
  recurrence~\citep{kalman1960}; recent neural versions such as Mamba handle long sequences
  efficiently~\citep{gu2024mamba}.
\gl{Token / tokenisation} The unit of input a sequence model reads / splitting data into such units.
\gl{Transformer} A sequence network built from attention layers; the architecture behind large
  language models~\citep{vaswani2017}.
\gl{Tree ensemble} Many decision trees combined, as in random forests or gradient
  boosting~\citep{hastie2009}.
\end{description}

\subsection{Probability and stochastic processes}
\begin{description}[leftmargin=1.2em,style=sameline,itemsep=1pt,font=\normalfont\bfseries]
\gl{Birth--death process} A counting process that moves up by one at a birth rate and down by one at a death rate that may depend on its current value~\citep{cst2010}.
\gl{Branching process} A model in which each individual independently produces a random number of offspring; a Hawkes process is one~\citep{hawkesoakes1974}.
\gl{Branching ratio} In a Hawkes process, the expected number of events directly triggered by one
  event (the spectral radius of the branching matrix in several dimensions); below one, every cascade dies out.
\gl{Cox--Ingersoll--Ross process} A continuous-time random process that is always positive and
  reverts toward a mean; used for interest rates and volatility.
\gl{Diffusive limit} The continuous random process that a jumpy process resembles when viewed at a
  large scale of time and size.
\gl{Ergodic} Having a unique long-run distribution, so that averages over a long history converge to fixed values.
\gl{Exogenous / endogenous} Caused from outside / triggered by earlier events of the same process.
\gl{Expectation--maximisation (EM)} An iterative fitting method for models with hidden quantities: guess the hidden quantities probabilistically from the current parameters, re-estimate the parameters, and repeat~\citep{veen2008em}.
\gl{First-passage time} The first time a random quantity reaches a given level.
\gl{Hawkes process} A self-exciting point process in which each event temporarily raises the rate
  of future events~\citep{hawkes1971,bacry2015}. Section~\ref{sec:hawkes}.
\gl{Heston model} A stochastic-volatility model in which variance follows a Cox--Ingersoll--Ross process.
\gl{Hurst exponent} A measure of long memory: $0.5$ for no memory, above $0.5$ for persistence.
\gl{Infimum ($\inf$)} The greatest lower bound of a set; for a set of times, the earliest one. By convention the infimum of an empty set is $+\infty$.
\gl{Intensity} The instantaneous expected rate of events.
\gl{Kalman filter} A recursive estimator of a hidden state that evolves linearly and is observed with noise~\citep{kalman1960}.
\gl{Laplace transform} For a waiting time $\tau$, the function $\varpi \mapsto \E[\mathrm{e}^{-\varpi\tau}]$; it turns sums of independent waiting times into products, which is why it is used for first-passage problems.
\gl{Markov} A process whose future depends only on its present state, not on its history.
\gl{Martingale} A process whose expected future value, given the present, equals its present value.
\gl{Point process} A random sequence of event times~\citep{daley2003}.
\gl{Queueing model} A probabilistic model of arrivals to and departures from queues.
\gl{Queue-reactive model} A model of the book in which the rate of each event type at a queue depends on the current queue sizes~\citep{huang2015}.
\gl{Stationary} Having statistical properties that do not change over time. For a Hawkes process,
  the cascade of triggered events does not explode.
\gl{Stochastic} Random.
\gl{Thinning} Simulating a point process by proposing candidate events at an upper-bound rate and keeping each with probability equal to the true intensity divided by the bound~\citep{ogata1981}.
\gl{Zumbach effect} Past volatility over long horizons predicts future short-horizon volatility better than the reverse; a time asymmetry of volatility~\citep{zumbach2009}.
\end{description}

\subsection{Systems and hardware}
\begin{description}[leftmargin=1.2em,style=sameline,itemsep=1pt,font=\normalfont\bfseries]
\gl{Actor model} A model of concurrency in which independent actors own their state and communicate
  only by messages~\citep{hewitt1973}.
\gl{ASIC} Application-specific integrated circuit: a chip designed for one task.
\gl{Chiplet} One of several small dies packaged together to act as a single processor.
\gl{Colocation} Placing trading servers in the exchange's data centre to minimise network distance.
\gl{FPGA} Field-programmable gate array: a chip whose logic can be reconfigured to implement a
  specific circuit; slower and larger than an ASIC for the same function but
  reprogrammable~\citep{kuon2007}.
\gl{GPU} Graphics processing unit: a processor with thousands of simple cores built for many arithmetic operations
  in parallel; a \emph{kernel} is one function launched on it, and a \emph{CUDA graph} records a sequence of kernels
  so it can be relaunched with less overhead.
\gl{High-level synthesis (HLS)} Compiling a program written in a language such as C++ into FPGA logic.
\gl{Jitter} Variation in latency from one event to the next.
\gl{Kernel bypass} Receiving and sending network packets directly from user space, skipping the
  operating system's network stack~\citep{dpdk}.
\gl{Multicast} One-to-many UDP delivery: the sender transmits once and every subscribed host
  receives the packet.
\gl{Operator mix} The kinds of computation a model needs, such as matrix multiplication or softmax.
\gl{PCIe} The bus that connects a processor to cards such as GPUs, FPGAs and network cards.
\gl{Queue-exact replay} A historical replay in which the simulated order sits in the reconstructed queue of recorded
  individual orders and is filled only when recorded executions reach an order that arrived after it; Section~\ref{sec:replay}.
\gl{Roofline model} A model of processor speed: attainable speed is the lower of peak arithmetic rate and memory
  bandwidth times operations per byte loaded~\citep{williams2009roofline}.
\gl{SmartNIC} A network card with programmable logic or processors that can compute on packets as they arrive.
\gl{STAC} An industry body that runs audited benchmarks of trading technology.
\gl{Tick-to-trade} The time from a market-data packet arriving to the resulting order leaving.
\gl{Transceiver} The circuit that converts between parallel data and the serial signal on a network link.
\gl{Trigger} Pre-computed order logic, often in hardware, that sends an order as soon as a market
  condition is detected.
\end{description}

\subsection{Economics}
\begin{description}[leftmargin=1.2em,style=sameline,itemsep=1pt,font=\normalfont\bfseries]
\gl{Agent-based model} A simulation of many artificial traders following rules and interacting.
\gl{Econometric / reduced-form model} A statistical model fitted to economic data that describes
  relationships without modelling traders' decisions.
\gl{Equilibrium model} A theoretical model in which prices result from rational traders' optimal
  strategies.
\gl{Glosten--Milgrom and Kyle models} Classic models in which market makers lose to better-informed
  traders and set prices or spreads to recover those losses~\citep{glosten1985,kyle1985}.
\gl{Latency arms race} Competition among fast traders to be first, and its economic
  costs~\citep{budish2015}.
\gl{Price formation} How trading and information are turned into the market price.
\end{description}
